% Preamble template laporan Milestone 2 IF4031 (mengacu pada gaya IF4031_M1_RMB.pdf)

\usepackage[a4paper,margin=2.54cm,headheight=15pt]{geometry}
\usepackage{iftex}

\ifPDFTeX
  \usepackage[utf8]{inputenc}
  \usepackage[T1]{fontenc}
  \usepackage{mathptmx}
  \usepackage[scaled=0.88]{helvet}
  \usepackage{courier}
\else
  \usepackage{fontspec}
  \IfFontExistsTF{Times New Roman}{
    \setmainfont{Times New Roman}
  }{
    \IfFontExistsTF{TeX Gyre Termes}{
      \setmainfont{TeX Gyre Termes}
    }{}
  }
  \IfFontExistsTF{Roboto Mono}{
    \setmonofont[Scale=0.88]{Roboto Mono}
  }{
    \IfFontExistsTF{Courier New}{
      \setmonofont[Scale=0.88]{Courier New}
    }{}
  }
\fi

\usepackage[indonesian]{babel}
\usepackage{graphicx}
\graphicspath{{figures/}{figures/proofs/}{../evidence/}}
\usepackage{float}
\usepackage{amsmath}
\usepackage{amssymb}
\usepackage{array}
\usepackage{booktabs}
\usepackage{tabularx}
\usepackage{longtable}
\usepackage{enumitem}
\usepackage{setspace}
\usepackage{parskip}
\usepackage{titlesec}
\usepackage{fancyhdr}
\usepackage{caption}
\usepackage{listings}
\usepackage{xcolor}
\usepackage{tikz}
\usetikzlibrary{positioning,arrows.meta,fit,backgrounds}
\usepackage{microtype}
\usepackage[normalem]{ulem}
\usepackage{hyperref}

\onehalfspacing
\emergencystretch=2.5em

% Warna selaras gaya dokumen M1 (#1155CC untuk tautan)
\definecolor{linkblue}{HTML}{1155CC}
\definecolor{codebg}{HTML}{F7F7F7}
\definecolor{codeborder}{HTML}{DDDDDD}
\definecolor{hintgray}{HTML}{555555}
\definecolor{hintbg}{HTML}{F9F9F9}

\hypersetup{
  colorlinks=true,
  linkcolor=black,
  urlcolor=linkblue,
  citecolor=linkblue,
  pdftitle={IF4031 Milestone 2 - RMB Distributed},
  pdfauthor={Kelompok 11: RMB}
}

% Format judul bab/subbab seperti IF4031_M1_RMB.pdf (14pt Bold + titik di belakang nomor)
\titleformat{\section}
  {\normalfont\fontsize{14}{17}\bfseries}{\thesection.}{0.6em}{}
\titleformat{\subsection}
  {\normalfont\fontsize{13}{16}\bfseries}{\thesubsection.}{0.6em}{}
\titleformat{\subsubsection}
  {\normalfont\fontsize{12}{15}\bfseries}{\thesubsubsection.}{0.6em}{}

\titlespacing*{\section}{0pt}{14pt}{6pt}
\titlespacing*{\subsection}{0pt}{11pt}{4pt}
\titlespacing*{\subsubsection}{0pt}{8pt}{3pt}

% Penomoran Gambar dan Tabel mengikuti hirarki subbab aktif (mis. Gambar 2.1.1, Tabel 3.1.3.1)
\makeatletter
\@addtoreset{figure}{section}
\@addtoreset{figure}{subsection}
\@addtoreset{figure}{subsubsection}
\@addtoreset{table}{section}
\@addtoreset{table}{subsection}
\@addtoreset{table}{subsubsection}
\renewcommand{\thefigure}{%
  \ifnum\value{subsubsection}>0 \thesubsubsection.\arabic{figure}%
  \else\ifnum\value{subsection}>0 \thesubsection.\arabic{figure}%
  \else\thesection.\arabic{figure}\fi\fi}
\renewcommand{\thetable}{%
  \ifnum\value{subsubsection}>0 \thesubsubsection.\arabic{table}%
  \else\ifnum\value{subsection}>0 \thesubsection.\arabic{table}%
  \else\thesection.\arabic{table}\fi\fi}
\makeatother

\captionsetup{font=small,labelfont=bf,justification=centering}
\captionsetup[table]{position=top,skip=4pt}
\captionsetup[figure]{position=bottom,skip=6pt}

% Header & footer
\pagestyle{fancy}
\fancyhf{}
\lhead{\small\textcolor{black}{IF4031 Arsitektur Aplikasi Terdistribusi}}
\rhead{\small\textcolor{black}{Kelompok 11: RMB}}
\cfoot{\small\thepage}
\renewcommand{\headrulewidth}{0.4pt}
\renewcommand{\footrulewidth}{0pt}

% Daftar dan kode sumber
\setlist[itemize]{leftmargin=1.5em,itemsep=2pt,topsep=4pt}
\setlist[enumerate]{leftmargin=1.5em,itemsep=2pt,topsep=4pt}

\lstset{
  basicstyle=\ttfamily\small,
  backgroundcolor=\color{codebg},
  frame=single,
  rulecolor=\color{codeborder},
  breaklines=true,
  columns=flexible,
  keepspaces=true,
  showstringspaces=false,
  tabsize=2
}

% Helper penanda isian kosong agar mudah dicari dan rapi di paragraf, list, maupun tabel
\newcommand{\isian}[1]{\textcolor{hintgray}{\textit{[ISIAN: #1]}}}

\newcommand{\blokisian}[1]{%
  \par\noindent\fcolorbox{codeborder}{hintbg}{%
    \parbox{\dimexpr\linewidth-2\fboxsep-2\fboxrule\relax}{%
      \small\textcolor{hintgray}{\textit{[ISIAN: #1]}}%
    }%
  }\par\vspace{4pt}%
}

\newcommand{\slotgambar}[2]{%
  \begin{figure}[H]
    \centering
    \fbox{\parbox[c][4.2cm][c]{0.85\linewidth}{\centering\small\textcolor{hintgray}{\textit{[Tempelkan gambar/screenshot: #1]}}}}
    % Ganti kotak di atas saat mengisi laporan dengan:
    % \includegraphics[width=0.92\linewidth]{nama-file.png}
    \caption{#2}
  \end{figure}
}

% Bukti pengujian: tangkapan layar hasil eksekusi pada klaster 3-VM (tinggi dibatasi agar tidak melebihi halaman)
\newcommand{\gambarbukti}[4][0.36]{%
  \begin{figure}[H]
    \centering
    \includegraphics[width=0.94\linewidth,height=#1\textheight,keepaspectratio]{#2}
    \caption{#3}\label{fig:#4}
  \end{figure}
}
